本研究以「預訓練—微調（Pre-train, Fine-tune）範式」與「Prompt Tuning 範式」為核心比較軸線，透過主實驗與一系列系統性消融實驗，對兩大訓練範式在語音深偽偵測任務上的跨域效能、架構極限及潛在改善空間進行了全面診斷。消融實驗的設計邏輯並非各自獨立，而是具有一貫的目的：\textbf{逐步探索 Prompt Tuning 範式能否在維持其高效訓練優勢的前提下，透過鬆綁前端離散化約束，逐步逼近預訓練—微調範式的偵測效能上限}。

本章具體總結如下：首先（\ref{sec:concl_paradigm_gap} 節）從範式層次對比主實驗結果，揭示兩種範式各自的跨域表現特性；其次（\ref{sec:concl_closing_gap} 節）系統性地梳理各消融實驗如何層層突破 Prompt Tuning 範式的離散化瓶頸；第三（\ref{sec:concl_mlaad} 節）針對 MLAAD 測試集異常的高準確率現象進行分析；第四（\ref{sec:concl_backend} 節）論證預訓練—微調範式在特徵聚合與後端強化設計上的重要性；最後（\ref{sec:concl_cost} 節）從訓練成本與參數效率的角度，對兩種範式進行綜合對比與觀點收斂。

\section{主實驗：兩大範式的跨域效能對比}
\label{sec:concl_paradigm_gap}

主實驗結果揭示，兩種訓練範式在跨域場景下均面臨效能衰退，但衰退的根源截然不同。

\textbf{預訓練—微調範式}（W2V-AASIST、MFA-Conformer）在同領域基準測試中展現出極高的偵測精準度（W2V-AASIST 於 ASV19 LA 達 0.27\% EER），其豐富的底層連續聲學特徵使其能精確捕捉偽造痕跡。然而，當測試場景轉移至存在未知編碼變異、通道失配或強烈環境噪音的跨域資料集時，所有模型的效能皆出現數十倍的大幅衰退。這顯示，預訓練—微調範式在單一資料集訓練時，容易對特定錄音環境或偽造演算法產生嚴重過擬合，尚未習得能夠應對真實世界多變場景的通用深偽辨別能力。

\textbf{Prompt Tuning 範式}（SpeechPrompt v2）則呈現出截然不同的困境。儘管透過極少量的 Prompt 參數仍能初步激發一定的分類能力，但即便在同領域測試中，其效能天花板即已遠低於預訓練—微調範式（ASV19 LA EER 18.50\% 對比 0.27\%）；在面對未曾見過的壓縮波形或各式未知偽造演算法時，其表現更是遠遠落後。

這一先天性的效能差距引出了一個核心問題：瓶頸究竟源於後端語言模型的學習能力不足，還是前端特徵處理過程本身的先天限制？此問題的答案，由下一節的系統性消融實驗逐步加以釐清。

\section{消融實驗：系統性逼近預訓練—微調範式的效能上限}
\label{sec:concl_closing_gap}

為了回答上述問題，本研究設計了一系列層層遞進的消融實驗，每一步均以前一步的發現為基礎，逐步鬆綁前端量化約束，系統性地探索 Prompt Tuning 範式的效能邊界。

\textbf{第一步：確認瓶頸位置（Prompt Tuning vs. Full Fine-Tuning）}。在確認介入方向之前，必須先釐清瓶頸究竟在後端語言模型還是前端離散化。實驗結果顯示，即使將整座 GSLM 語言模型解凍進行全參數微調，效能提升極為有限（ASV19 LA EER 僅從 18.50\% 降至 17.23\%），而在難度更高的 ASV21 DF 上甚至略微退步。此結果直接排除了「後端語言模型能力不足」的假設，將介入矛頭明確指向前端離散化過程，為後續消融實驗確立了正確的研究方向。

\textbf{第二步：部分保留連續資訊（軟性離散特徵）}。以 Softmax 加權代替硬性最近鄰配對，使連續特徵在映射時保留各群集中心的相對距離資訊，將 ASV19 LA 的 EER 從 18.50\% 改善至 15.90\%。改善幅度雖然有限，但初步驗證了「保留更多連續資訊確實有助於提升 Prompt Tuning 範式偵測能力」的假設，並顯示靜態群集中心佈局仍是制約效能的次要瓶頸。

\textbf{第三步：釋放群集中心約束（可學習群集中心與直接線性映射）}。此為本研究最關鍵的消融實驗。群集中心殘差微調（Centroid Residual Tuning）使 EER 大幅降至 5.53\%；完全移除量化器的直接線性映射則進一步達到 3.84\%，在同域測試集（ASV19 LA）上已大幅逼近預訓練—微調範式（W2V-AASIST 0.27\%、MFA-Conformer 1.12\%）的效能水準；跨域測試集上的改善同樣一致（ASV21 DF 由基準 28.94\% 大幅降至 14.44\%），惟與 W2V-AASIST 在 ASV21 DF 上的 5.57\% 相比仍有可觀差距。

三步實驗共同呈現出一條清晰的遞進曲線：
\begin{center}
硬性量化（18.50\%）$\rightarrow$ 軟性離散（15.90\%）$\rightarrow$ 群集中心殘差微調（5.53\%）$\rightarrow$ 直接線性映射（3.84\%）
\end{center}
隨著前端量化約束的逐步鬆綁，Prompt Tuning 範式的效能上限持續突破，且改善幅度隨約束鬆綁程度而顯著增大。此結果有力地實證了先前的假設：深偽音訊與真實語音之間的本質差異，往往存在於極其微小且難以察覺的連續聲學特徵變化或高頻響應扭曲中，而硬性量化到有限集合離散單元的過程，實質上等同於一個極端不可逆的低通濾波器，直接將這類用以區分真偽的關鍵微小瑕疵視為「雜訊」加以抹除，使系統從源頭便喪失了建立穩健防禦決策所需的關鍵高頻與細節特徵。\textbf{硬性離散化造成的量化資訊損失，正是 Prompt Tuning 範式效能遠落後於預訓練—微調範式的最根本原因}；而若能從前端根本上消除此約束，Prompt Tuning 範式在同域條件下已具備大幅逼近連續特徵架構效能的潛力；跨域泛化的進一步突破，則仍有賴語言模型骨幹架構層面的根本性改進。

\section{MLAAD 測試集的異常高準確率分析}
\label{sec:concl_mlaad}
在跨領域測試中，我們觀察到一個極度反常的現象：在多數測試集中表現最差、幾乎不具備跨域能力（EER 極高）的 SpeechPrompt v2，卻在 MLAADv2 與 MLAADv8 上達到將近 100\% 的完美準確率；反而是在其他標準與真實場景測試集上最為穩健的 W2V-AASIST，在此資料集上相對表現不佳（準確率約七至八成）。

針對此一反轉，我們認為並非 SpeechPrompt 具備某種未知的強大泛化性，而是反映了 MLAAD 資料集本身存在潛在的特性落差。具體而言：
\begin{itemize}
    \item \textbf{資料集潛在的捷徑 (Shortcut) 特徵}：MLAAD 的假音源或採樣方式可能帶有某種極為固定的背景噪音、特定編碼偽影或特殊的頻譜破壞，而這種單一破壞剛好對應到離散化模型中的特定單元群。使得模型無須學習「深偽本質」，只需退化成一個簡單的「特定離散單元或噪聲分類器」即可完美分類。而恰好 SpeechPrompt 的語言模型架構本身便非常擅長捕捉這類離散序列的出現頻率。
    \item \textbf{連續特徵對未知極端變異的反效果}：相較之下，W2V-AASIST 依賴更細緻的連續特徵學習，當 MLAAD 中存在與 ASV19 或 InTheWild 截然不同的強烈聲學特徵變異（如極端殘響或極差音質）時，過大的變異可能反而跨越了判別邊界，導致原本仰賴微小連續分佈特徵的分類網路發生混淆。
\end{itemize}
此現象進一步提示：在評估深偽偵測系統時，不同資料集的分佈特性可能對兩種範式產生截然相反的影響，單一資料集上的高準確率並不等同於真實場景下的強健泛化能力。

\section{預訓練—微調範式之擴容局限與後端特徵融合的必要性}
\label{sec:concl_backend}
對於預訓練—微調範式本身，消融實驗同樣帶來了重要啟示。MFA-Conformer 的模型規模擴容實驗證實，即使將前端從 Conformer-Small（13.1M 參數）大幅擴增至 Conformer-Large（121.5M 參數），在 ASV21 DF 的 EER 依然停滯在 13\%--14\% 附近，呈現明顯的效能飽和狀態。這說明在預訓練—微調範式中，單純堆疊更大的前端模型並非突破跨域瓶頸的有效路徑。

真正帶來效能突破的，是改變後端的特徵融合架構。當我們保持前端特徵不變，將後端從簡單的線性分類器替換為能進行多層次圖結構融合的 AASIST 架構，並引入所有層次的 Conformer 特徵（All layer），ASV21 DF 的 EER 從 16.00\% 大幅改善至 12.71\%。這說明在預訓練—微調範式中，\textbf{前端多層次特徵的充分聚合能力}與\textbf{後端的複雜分類決策}兩者缺一不可，僅有強大的特徵提取前端，而缺乏能充分整合多尺度聲學資訊的後端架構，依然無法釋放前端模型的完整潛能。

\section{訓練成本、參數效率與兩大範式的綜合觀點}
\label{sec:concl_cost}
除偵測準確率外，兩種範式在訓練成本與部署效率上呈現出截然懸殊的差距，此差距本身即是重要的實驗發現。

預訓練—微調範式中，W2V-AASIST 由於需對 317M 參數的前端模型進行全域微調且直接處理原始波形，計算負擔極重，在設定 100 Epoch 的情況下，單次完整訓練耗時達 \textbf{24 至 36 小時}；MFA-Conformer（14.6M）得益於頻譜特徵的維度縮減，訓練速度較快，在固定 50 Epoch 的設定下，完整訓練約需 \textbf{2 小時}。Prompt Tuning 範式的 SpeechPrompt v2 則以凍結主幹、僅更新 0.13M 個 Prompt 參數的方式，由於輸入為預先計算好的離散 Token 且模型收斂極快，將完整訓練時間壓縮至約 \textbf{300 秒（5 分鐘）}，實現了近 400 倍的訓練時間優勢。

綜合本研究所有實驗發現，兩種範式呈現出高度互補的特性：預訓練—微調範式具備較高的偵測效能上限，但訓練成本高昂，且存在效能飽和與跨域泛化的結構性挑戰；Prompt Tuning 範式訓練效率極高，然而當前以離散 Token 為輸入前提的語音語言模型架構，使其受制於前端量化瓶頸，難以充分發揮潛力。消融實驗的遞進結果已清楚指向：若未來能以\textbf{原生支援連續聲學特徵輸入的新一代語音語言模型}替換當前以離散 Token 為前提的 GSLM 骨幹，使 Prompt Tuning 範式得以在不損失細粒度聲學資訊的前提下運作，兩種範式各自的優勢將有望在同一架構中融合——兼具預訓練—微調範式的高效能上限，以及 Prompt Tuning 範式的極低訓練成本，從而為未來的深偽防禦系統提供更具實用潛力的解決路徑。
